% Part: history
% Chapter: biographies
% Section: rozsa-peter

\documentclass[../../../include/open-logic-section]{subfiles}

\begin{document}

\olfileid{his}{bio}{pet}

\olsection{রোজা পেতের}

\olphoto{peter-rozsa}{রোজা পেতের}

রোজা পেতেরের জন্মনাম ছিল রোজা পোলিৎসার। ১৯০৫
সালের ১৭ ফেব্রুয়ারি হাঙ্গেরির বুদাপেস্টে তাঁর জন্ম।
পুনরাবৃত্ত অপেক্ষক নিয়ে কাজের জন্যই তিনি সবচেয়ে
পরিচিত; পুনরাবৃত্তি তত্ত্বের ক্ষেত্রটি গড়ে উঠতে এই
কাজ অপরিহার্য ছিল।

পেতেরের বেড়ে ওঠার সময়ে রাজনৈতিক পরিস্থিতি ছিল
কঠিন—তাঁর কৈশোরেই প্রথম বিশ্বযুদ্ধ চলছিল। তবু
তিনি বুদাপেস্টের সচ্ছল মারিয়া তেরেজিয়া বালিকা
বিদ্যালয়ে পড়তে পেরেছিলেন; সেখান থেকে ১৯২২
সালে পড়া শেষ করেন। তারপর বুদাপেস্টের পাজমানি
পেতের বিশ্ববিদ্যালয়ে (পরে যার নাম হয় লোরান্ড
এয়োটভোশ বিশ্ববিদ্যালয়) পড়েন। বাবার জোরাজুরিতে
প্রথমে রসায়ন নিয়ে পড়া শুরু করলেও পরে গণিতে
চলে আসেন এবং ১৯২৭ সালে স্নাতক হন। উচ্চবিদ্যালয়ে
গণিত পড়ানোর যোগ্যতা থাকলেও মহামন্দায় বিশ্ব
অর্থনীতি বিপর্যস্ত হওয়ায় তখন কর্মসংস্থানের
পরিস্থিতি ছিল ভয়াবহ। এই সময়ে পেতের গৃহশিক্ষকতা
ও ব্যক্তিগতভাবে গণিত পড়ানোর নানা কাজ করেন।
শেষ পর্যন্ত তিনি গণিতে উচ্চতর পড়াশোনার জন্য
বিশ্ববিদ্যালয়ে ফেরেন। গোড়ায় সংখ্যাতত্ত্ব নিয়ে
কাজ করার পরিকল্পনা ছিল; কিন্তু তাঁর ফলগুলি
আগেই অন্যরা প্রমাণ করেছেন জানতে পেরে প্রায়
গণিতই ছেড়ে দিচ্ছিলেন। পরে গ্যোডেলের অসম্পূর্ণতা
উপপাদ্যগুলি নিয়ে কাজ করতে তাঁকে উৎসাহ দেওয়া
হয়; অজান্তেই তিনি গ্যোডেলের বেশ কয়েকটি ফল
ভিন্ন উপায়ে প্রমাণ করেন। তাতে তাঁর আত্মবিশ্বাস
ফিরে আসে। ডাভিড হিলবার্টের ভিত্তি-সংক্রান্ত
কর্মসূচি থেকে অনুপ্রেরণা নিয়ে পেতের পুনরাবৃত্তি
তত্ত্বের উপর তাঁর প্রথম প্রবন্ধগুলি লেখেন। ১৯৩৫
সালে তিনি পিএইচডি লাভ করেন; ১৯৩৭ সালে
\emph{Journal of Symbolic Logic}-এর সম্পাদক হন।

পুনরাবৃত্ত অপেক্ষক তত্ত্বের প্রতিষ্ঠায় পেতেরের
প্রথমদিকের প্রবন্ধগুলির অবদান ব্যাপকভাবে স্বীকৃত।
\cite{Peter1935a}-এ তিনি বিভিন্ন ধরনের
পুনরাবৃত্তির পারস্পরিক সম্পর্ক অনুসন্ধান করেন।
\cite{Peter1935b}-এ তিনি দেখান যে পুনরাবৃত্তভাবে
সংজ্ঞায়িত একটি নির্দিষ্ট অপেক্ষক আদিম পুনরাবৃত্ত
নয়। এতে ভিলহেল্ম অ্যাকারমানের আগের একটি
ফল সরল হয়। পেতেরের সরলীকৃত অপেক্ষকটিকেই
এখন প্রায়ই অ্যাকারমান অপেক্ষক বলা হয়—কখনো
আরও যথাযথভাবে অ্যাকারমান--পেতের অপেক্ষক।
পুনরাবৃত্ত অপেক্ষক তত্ত্বের প্রথম বইটিও তিনি লেখেন
\citep{Peter1951}।

তাঁর কাজের গুরুত্ব ও প্রভাব সত্ত্বেও ১৯৪৫ সালের
আগে পেতের পূর্ণকালীন শিক্ষকতার পদ পাননি।
দ্বিতীয় বিশ্বযুদ্ধকালে হাঙ্গেরিতে নাৎসি দখলদারির
সময় ইহুদিবিদ্বেষী আইনে তাঁকে পড়াতে দেওয়া হয়নি।
১৯৪৪ সালে সরকার বুদাপেস্টে একটি ইহুদি ঘেটো
তৈরি করে; সশস্ত্র পাহারায় সেটিকে শহরের বাকি
অংশ থেকে বিচ্ছিন্ন রাখা হয়। ১৯৪৫ সালে মুক্তি
পাওয়ার আগে পর্যন্ত পেতেরকে সেখানে থাকতে
বাধ্য করা হয়। পরে তিনি বুদাপেস্ট শিক্ষক-প্রশিক্ষণ
মহাবিদ্যালয়ে এবং ১৯৫৫ সাল থেকে এয়োটভোশ
লোরান্ড বিশ্ববিদ্যালয়ে পড়ান। হাঙ্গেরির প্রথম
নারী গণিতবিদ হিসেবে তিনি গণিতে অ্যাকাডেমিক
ডক্টর পদমর্যাদা পান এবং হাঙ্গেরীয় বিজ্ঞান
অ্যাকাডেমির সদস্য নির্বাচিত হন।
% BN-SRC-850 TARGET-CORRECTION-BEGIN
\par\noindent\textbf{পাঠ-সংশোধন।} মূলপাঠের “প্রথম নারী” সংশোধিত: এর্জেবেত আন্দিচ
১৯৪৯ সালে এবং পেতের ১৯৭৩ সালে সদস্য নির্বাচিত হন। দেখো
\href{https://akademikus.mtak.hu/adatlap/andics-erzsebet/}{আন্দিচের নথি} ও
\href{https://akademikus.mtak.hu/adatlap/peter-rozsa/}{পেতেরের নথি}।
% BN-SRC-850 TARGET-CORRECTION-END

পেতের ছিলেন গণিতের এক নিবেদিত শিক্ষক।
শুধু বক্তৃতা দেওয়ার চেয়ে শিক্ষার্থীদের সঙ্গে
গাণিতিক সমস্যার প্রকৃতি ও সৌন্দর্য অনুসন্ধান
করতেই তিনি বেশি পছন্দ করতেন। তাই
শিক্ষার্থীরা স্নেহ করে তাঁকে “রোজা মাসি”
(Aunt Rosa) বলতেন। ১৯৭৭ সালে ৭১ বছর
বয়সে তাঁর মৃত্যু হয়।

\begin{reading}
আরও জীবনীমূলক রচনার জন্য \citep{Oconnor2014}
ও \citep{Andrasfai1986} দেখো। পেতেরের একটি
সংক্ষিপ্ত সাক্ষাৎকার নিয়েছিলেন
\citet{Tamassy1994}। গণিত নিয়ে আনন্দদায়ক
পাঠের জন্য পেতেরের \emph{Playing With Infinity}
বইটি দেখো \citep{Peter2010}।
\end{reading}

\end{document}
